\section{Methodology}
\label{sec:method}

We present and analyze \textbf{Parameter-Free Task-Space Projection (PFSR)}, a remarkably simple, elegant, and training-free dynamic ensembling framework. Guided by our \textit{Minimalist} philosophy, PFSR introduces zero trainable parameters and requires zero training epochs or calibration splits. Instead, it extracts task-representative centroids from pre-trained specialists to perform direct online coordinate projection. 

To investigate whether we can decouple task coordinate axes and eliminate cross-talk under task overlap, we also propose and analyze an advanced extension: \textbf{L{\"o}wdin-Orthogonalized Task-Space Projection (OTSP)}. OTSP computes an optimal, order-invariant orthonormal task coordinate basis $\{q_k\}_{k=1}^K$ in closed form using the symmetric inverse square root of the Gram overlap matrix, directly from the frozen classification weights of $K$ pre-trained specialists. The methodology proceeds in five clean mathematical steps.

\subsection{Step 1: Extract Task Centroids via Singular Value Decomposition}
To obtain a representative directional vector for the semantic subspace of Expert $k$, we must extract a task-specific centroid vector $v_k \in \mathbb{R}^D$ from its weight matrix $W_k \in \mathbb{R}^{C \times D}$.
In standard trained networks, final classifier class prototypes (the rows of $W_k$) are symmetrically distributed around the origin (sum-to-zero) to maximize classification margins. Consequently, taking a simple average of all class prototypes $W_{k,c}$ would yield a vector extremely close to $\mathbf{0}$, where normalization merely amplifies random numerical noise.
To resolve this cancellation, we extract the task centroid using \textbf{Singular Value Decomposition (SVD)}. We perform SVD on the expert's weight matrix:
\begin{equation}
W_k = U_k \Sigma_k V_k^T
\end{equation}
where $V_k \in \mathbb{R}^{D \times D}$ is an orthogonal matrix whose columns represent the principal components of the prototype subspace. We define the raw task centroid $v_k$ as the top right-singular vector (the first column of $V_k$, or equivalently, the first row of $V_k^T$):
\begin{equation}
v_k = V_{k, 1}
\end{equation}
This captures the direction of maximum variance of the class prototypes, which is perfectly stable and completely immune to sum-to-zero cancellation. We normalize $v_k$ to unit-norm to obtain the raw task direction $\bar{v}_k$:
\begin{equation}
\bar{v}_k = \frac{v_k}{\|v_k\|_2}
\end{equation}

\textbf{Alternative Centroid Formulations (Weights vs. Activations):} While our primary formulation extracts centroids directly from the final classifier weights $W_k$, this assumes that the explicit classification heads are accessible. In settings where classifier weights are not directly available (e.g., intermediate layers, or models merged without explicit head weights), alternative data-dependent centroid formulations can be utilized. Specifically, given a tiny calibration set of sample activations $H_k \in \mathbb{R}^{N_{\text{cal}} \times D}$ for task $k$, one can compute the centroid as: (1) the mean-normalized representation $\bar{h}_k = \frac{1}{N_{\text{cal}}} \sum_{i} H_{k, i}$, which is robust to prototype cancellation because activations do not exhibit sum-to-zero symmetry; (2) the first principal component derived from SVD on the activation matrix $H_k$; or (3) the principal cluster centroid using $K$-Means clustering on the representations.

These alternative formulations exhibit distinct practical and computational trade-offs:
\begin{itemize}
    \item \textbf{SVD on Classifier Weights (Primary):} This is 100\% data-free, requiring zero calibration samples or forward-pass activations. It relies solely on the static parameters of the pre-trained classification layer. It is computationally instantaneous ($O(K C D)$ where $C$ is class cardinality) and does not incur any offline data-storage or privacy overhead. However, it requires direct parameter access to the classification heads, making it less applicable to intermediate layers.
    \item \textbf{Mean/SVD on Activations:} These options require a small representative dataset of sample activations (e.g., a 64-sample calibration split). Although this introduces a minor data-dependency, activation-based centroids capture the actual, contextualized distribution of the model's feature space. This makes them highly robust to out-of-distribution drift, and more importantly, allows task-space projection to be applied at any arbitrary intermediate hidden layer in the model, even where explicit classification heads do not exist.
    \item \textbf{K-Means Clustering on Activations:} Running $K$-Means on representation activations identifies multiple sub-clusters or local semantic centroids within a task's feature manifold. While computationally more expensive than simple averaging or SVD ($O(I K N_{\text{cal}} D)$ where $I$ is iteration count), it successfully captures multi-modal task distributions, allowing for highly flexible routing in complex multi-task environments where a single linear centroid is insufficient to represent the entire task space.
\end{itemize}

\subsection{Step 2: Compute Overlap (Gram) Matrix}
Independently trained specialists often share overlapping feature structures, leading to directional correlation between their raw task directions. To quantify this overlap, we construct the symmetric Gram matrix $S \in \mathbb{R}^{K \times K}$ representing the pairwise cosine similarities of the task directions:
\begin{equation}
S_{ij} = \bar{v}_i \cdot \bar{v}_j
\end{equation}

\subsection{Step 3: L{\"o}wdin Symmetric Orthogonalization}
To decouple the task coordinate spaces while preserving mathematical symmetry and order-invariance, we apply L{\"o}wdin Symmetric Orthogonalization~\cite{lowdin1950non}. This operation finds the set of orthonormal vectors $\{q_k\}_{k=1}^K$ that minimizes the least-squares distance to the original normalized centroids:
\begin{equation}
\min_{\{q_k\}} \sum_{k=1}^K \|q_k - \bar{v}_k\|_2^2 \quad \text{subject to } q_i \cdot q_j = \delta_{ij}
\end{equation}
Because the Gram matrix $S$ is symmetric positive semidefinite, we compute its eigendecomposition:
\begin{equation}
S = U \Lambda U^T
\end{equation}
where $U \in \mathbb{R}^{K \times K}$ is an orthogonal matrix of eigenvectors, and $\Lambda = \text{diag}(\lambda_1, \dots, \lambda_K)$ represents the eigenvalues. To eliminate task correlation while maintaining order-invariance, we compute the symmetric inverse square root of $S$:
\begin{equation}
\begin{aligned}
S^{-1/2} &= U \Lambda^{-1/2} U^T \\
&= U \text{diag}\left(\frac{1}{\sqrt{\lambda_1 + \epsilon}}, \dots, \frac{1}{\sqrt{\lambda_K + \epsilon}}\right) U^T
\end{aligned}
\end{equation}
where $\epsilon = 10^{-6}$ is a numerical stabilizer. The symmetric orthonormalized task directions $\{q_1, \dots, q_K\} \in \mathbb{R}^D$ are computed as a linear combination of the original centroids:
\begin{equation}
q_k = \sum_{j=1}^K (S^{-1/2})_{kj} \bar{v}_j
\end{equation}
By construction, the resulting basis is perfectly orthonormal: $q_i \cdot q_j = \delta_{ij}$ for all $i, j \in \{1, \dots, K\}$. This offline transformation takes less than a millisecond and is executed completely data-free.

\subsection{Step 4: Task Coordinate Absolute Projection}
During dynamic online inference, the input batch $X$ is passed through the frozen pre-trained base model backbone to extract the penultimate feature representations. Let $z_b \in \mathbb{R}^D$ be the penultimate representation of sample $b$ in the batch. We normalize the feature representation to unit-norm to ensure scale-invariance:
\begin{equation}
\tilde{z}_b = \frac{z_b}{\|z_b\|_2}
\end{equation}
Because class prototypes point in opposite directions within the task's subspace (symmetrically distributed around the origin), a simple linear projection will yield positive projections for some classes and negative projections for others, resulting in systematic routing failures within the same task. To resolve this, we take the \textbf{absolute value} of the projection coordinates.

For our primary, minimalist \textbf{PFSR} router, the coordinates $u_b = [u_{1, b}, \dots, u_{K, b}]^T \in \mathbb{R}^K$ are computed by taking the absolute projection onto the raw normalized centroids:
\begin{equation}
u_{k, b} = |\bar{v}_k \cdot \tilde{z}_b|
\end{equation}
For the \textbf{OTSP} extension, we compute the coordinate absolute projection $u'_b = [u'_{1, b}, \dots, u'_{K, b}]^T \in \mathbb{R}^K$ by projecting onto our orthonormal task basis $Q \in \mathbb{R}^{K \times D}$ and taking the absolute value:
\begin{equation}
u'_{k, b} = |q_k \cdot \tilde{z}_b|
\end{equation}
Taking the absolute value ensures that both positive and negative directions within a task's subspace yield high projection scores, matching both halves of the symmetric class distribution, whereas $u_{k,b}$ measures raw directional similarity.

\subsection{Step 5: Temperature-Scaled Softmax Gating}
The sample-wise dynamic merging coefficients $\alpha_b \in \mathbb{R}^K$ are computed by applying a temperature-scaled Softmax function over the projection coordinates ($u_{k, b}$ for PFSR, or $u'_{k, b}$ for OTSP):
\begin{equation}
\alpha_{k, b} = \frac{\exp(\hat{u}_{k, b} / \tau)}{\sum_{j=1}^K \exp(\hat{u}_{j, b} / \tau)}
\end{equation}
where $\hat{u}_{k, b}$ represents $u_{k, b}$ for PFSR or $u'_{k, b}$ for OTSP, and $\tau > 0$ is a static scaling temperature. In practice, we set $\tau = 0.001$, which acts as a sharp selector to route samples with near-discrete task specialization. The resulting coefficients $\alpha_b$ satisfy the ensembling simplex constraints ($\sum_k \alpha_{k,b} = 1$ and $\alpha_{k,b} \ge 0$). These coefficients are then used to merge the LoRA adapters on the fly, performing a single forward pass over the merged parameters for the corresponding sequence.

We explicitly clarify that while PFSR and OTSP are strictly \textit{parameter-free} in terms of \textit{trainable parameters}---requiring exactly zero weight matrices, offline optimization runs, or backpropagation calculations---the dynamic gating remains sensitive to the choice of the temperature hyperparameter $\tau$. This temperature operates as a continuous dial controlling the fundamental ensembling trade-off. Under sharp gating ($\tau \to 0$), the Softmax converges to a hard argmax selector, offering robust \textit{noise isolation} and specialist protection but suffering from a \textit{hard gating penalty} by sacrificing prediction-averaging benefits under active task overlap. Conversely, under smooth gating ($\tau \to \infty$), the routing weights approach a static uniform ensembling state. To achieve a completely self-contained online deployment with zero manual calibration, we propose a simple, self-calibrated temperature scheduling heuristic: setting a sample-wise dynamic temperature $\tau_b = \gamma \cdot \text{std}_k(u_{k, b})$, where $\text{std}_k(u_{k,b})$ is the standard deviation of the projection coordinates across the experts for sample $b$, and $\gamma > 0$ is a constant scaling multiplier. This online self-calibration adapts the routing sharpness dynamically to sample ambiguity (higher temperature for high-uncertainty overlap regions, lower temperature for clear specialization) with zero calibration split overhead, ensuring robust performance across varying registries.

\subsection{The Noise Amplification Penalty in Orthonormal Projection}
\label{sec:noise_amplification}
While L{\"o}wdin Symmetric Orthogonalization is mathematically elegant, a deeper geometric and numerical analysis reveals a fundamental limitation of the orthonormal projection paradigm under heavy task overlap. This issue is analogous to the problem of \textit{multicollinearity} in classical linear regression.

Let the original normalized task centroids $\{\bar{v}_k\}$ share a high pairwise cosine similarity (e.g., $\bar{v}_1 \cdot \bar{v}_2 \to 1$). In this regime, the Gram matrix $S$ becomes near-singular, and its smallest eigenvalue $\lambda_{\min}$ approaches zero. When we compute the symmetric inverse square root matrix $S^{-1/2} = U \Lambda^{-1/2} U^T$, the scaling factors $1/\sqrt{\lambda_j + \epsilon}$ for the eigenvectors corresponding to small eigenvalues become extremely large.

Consequently, to achieve mathematical orthonormality, the resulting basis vectors $q_k = \sum_j (S^{-1/2})_{kj} \bar{v}_j$ must use large coefficients with alternating signs to cancel out the directional overlap. For instance, in a two-task setting with high correlation, we have $q_1 \approx a \bar{v}_1 - b \bar{v}_2$ where $a, b \gg 1$. 

During online dynamic inference, let the normalized feature representation of sample $b$ be modeled as a true representation $\tilde{z}_b^*$ corrupted by isotropic representation noise $\eta_b$:
\begin{equation}
\tilde{z}_b = \tilde{z}_b^* + \eta_b
\end{equation}
When we project this noisy representation onto our orthonormal task basis, the coordinate $u'_{k,b}$ becomes:
\begin{equation}
u'_{k,b} = q_k \cdot \tilde{z}_b = q_k \cdot \tilde{z}_b^* + q_k \cdot \eta_b
\end{equation}
Because the basis vector $q_k$ contains large scaling coefficients from $S^{-1/2}$, the noise projection term $q_k \cdot \eta_b$ undergoes severe \textbf{noise amplification}. The variance of the projection coordinate is scaled by the squared norm of the inverse square root coefficients:
\begin{equation}
\text{Var}(q_k \cdot \eta_b) = \sigma^2 \|q_k\|_2^2 = \sigma^2 (S^{-1})_{kk}
\end{equation}
where $(S^{-1})_{kk} \gg 1$ under heavy overlap, and $\sigma^2$ is the variance of the isotropic feature noise.

This noise amplification destabilizes the coordinate values and routing decisions compared to the raw unorthogonalized projection (PFSR). Under PFSR, the projection coordinates are simply $u_{k,b} = \bar{v}_k \cdot \tilde{z}_b$, and because $\|\bar{v}_k\|_2 = 1$, the noise projection term $\bar{v}_k \cdot \eta_b$ has stable variance $\sigma^2$ with zero amplification. Thus, we identify a fundamental minimalist trade-off: orthogonalizing the task axes successfully decorrelates the offline templates, but under high overlap and active representation noise, it amplifies online coordinate variance, leading to potential routing degradation.

\subsection{Mathematical Equivalence under Symmetric Task Overlap}
\label{sec:symmetric_equivalence}
To formally address the relationship between OTSP and the unorthogonalized PFSR, we analyze the routing decisions of both methods under symmetric task overlap. 

Consider a two-task setting ($K=2$) with symmetric centroid similarity $s = \bar{v}_1 \cdot \bar{v}_2 \in [0, 1)$. The Gram overlap matrix is given by:
\begin{equation}
S = \begin{pmatrix} 1 & s \\ s & 1 \end{pmatrix}
\end{equation}
The eigenvalues of $S$ are $\lambda_1 = 1+s$ and $\lambda_2 = 1-s$, with corresponding orthonormal eigenvectors:
\begin{equation}
w_1 = \frac{1}{\sqrt{2}} [1, 1]^T, \quad w_2 = \frac{1}{\sqrt{2}} [1, -1]^T
\end{equation}
Computing the symmetric inverse square root $S^{-1/2} = U \Lambda^{-1/2} U^T$ yields:
\begin{equation}
S^{-1/2} = \begin{pmatrix} a & b \\ b & a \end{pmatrix}
\end{equation}
where the diagonal and off-diagonal elements are:
\begin{equation}
\begin{aligned}
a &= \frac{1}{2}\left( \frac{1}{\sqrt{1+s}} + \frac{1}{\sqrt{1-s}} \right), \\
b &= \frac{1}{2}\left( \frac{1}{\sqrt{1+s}} - \frac{1}{\sqrt{1-s}} \right)
\end{aligned}
\end{equation}
The L{\"o}wdin orthonormal basis vectors $\{q_1, q_2\}$ are computed as:
\begin{equation}
q_1 = a \bar{v}_1 + b \bar{v}_2, \quad q_2 = b \bar{v}_1 + a \bar{v}_2
\end{equation}
To determine the routing decision, we examine the coordinate difference $u'_1 - u'_2$ for any sample $\tilde{z}_b$:
\begin{equation}
u'_1 - u'_2 = (q_1 - q_2) \cdot \tilde{z}_b = (a - b) (\bar{v}_1 - \bar{v}_2) \cdot \tilde{z}_b
\end{equation}
Substituting the definitions of $a$ and $b$, the difference $a - b$ simplifies to:
\begin{equation}
a - b = \frac{1}{\sqrt{1-s}}
\end{equation}
Therefore, the OTSP coordinate difference is directly proportional to the raw PFSR coordinate difference ($u_1 - u_2$):
\begin{equation}
u'_1 - u'_2 = \frac{1}{\sqrt{1-s}} (u_1 - u_2)
\end{equation}
Because the scaling factor $\frac{1}{\sqrt{1-s}} > 0$ is strictly positive for all $s \in [0, 1)$, we have:
\begin{equation}
\text{sign}(u'_1 - u'_2) = \text{sign}(u_1 - u_2)
\end{equation}
This mathematically guarantees that the argmax of OTSP and PFSR are \textbf{identical for every single sample}, regardless of the presence or magnitude of representation noise $\eta_b$. Under a near-hard gating temperature (e.g., $\tau = 0.001$), the Softmax gating coefficients $\alpha_b$ are completely determined by the argmax, forcing OTSP and PFSR to make \textbf{exactly identical routing and merging decisions}.

In asymmetric settings, however, the Gram matrix $S$ does not possess a constant off-diagonal algebraic structure. Consequently, the inverse square root $S^{-1/2}$ does not scale the coordinate differences uniformly, allowing OTSP and PFSR to make different routing decisions. As shown empirically in Section~\ref{sec:experiments}, under asymmetric overlap and moderate noise, OTSP's orthogonalization suffers from the Noise Amplification Penalty, slightly degrading its performance relative to the raw, unorthogonalized PFSR.

\subsection{Signal-to-Noise Ratio (SNR) Equivalence}
\label{sec:snr_equivalence}
To understand why L{\"o}wdin orthogonalization fails to provide any routing accuracy gains even in noisy settings under task overlap, we derive the exact Signal-to-Noise Ratio (SNR) of the coordinate difference under both methods.

Let a feature representation of a sample belonging to Task 1 be modeled as $z_b = \bar{v}_1 + \eta_b$, where $\eta_b$ is isotropic representation noise with variance $\sigma^2$ (i.e., $\mathbb{E}[\eta_b \eta_b^T] = \sigma^2 I_D$).
Under unorthogonalized projection (PFSR), the coordinate difference is $u_1 - u_2 = (\bar{v}_1 - \bar{v}_2) \cdot z_b$. The clean (noise-free) coordinate difference is:
\begin{equation}
\Delta u^* = (\bar{v}_1 - \bar{v}_2) \cdot \bar{v}_1 = 1 - s
\end{equation}
The variance of the noise projection on the difference axis is:
\begin{equation}
\text{Var}((\bar{v}_1 - \bar{v}_2) \cdot \eta_b) = \sigma^2 \|\bar{v}_1 - \bar{v}_2\|_2^2 = \sigma^2 (2 - 2s)
\end{equation}
Taking the ratio of the clean margin to the noise standard deviation, the Signal-to-Noise Ratio under PFSR is:
\begin{equation}
\text{SNR}_{\text{PFSR}} = \frac{1-s}{\sigma \sqrt{2(1-s)}} = \frac{\sqrt{1-s}}{\sigma \sqrt{2}}
\end{equation}

Under our proposed L{\"o}wdin-orthogonalized projection (OTSP), the coordinate difference is $(q_1 - q_2) \cdot z_b$. The clean coordinate difference (expanded margin) is:
\begin{equation}
\begin{aligned}
\Delta u'^* &= (q_1 - q_2) \cdot \bar{v}_1 \\
&= (a - b)(\bar{v}_1 - \bar{v}_2) \cdot \bar{v}_1 \\
&= \frac{1-s}{\sqrt{1-s}} = \sqrt{1-s}
\end{aligned}
\end{equation}
The variance of the noise projection under OTSP is:
\begin{equation}
\text{Var}((q_1 - q_2) \cdot \eta_b) = \sigma^2 \|q_1 - q_2\|_2^2
\end{equation}
Because the basis vectors $\{q_k\}$ are perfectly orthonormal, we have $q_1 \cdot q_2 = 0$ and $\|q_1\|_2 = \|q_2\|_2 = 1$. Therefore, $\|q_1 - q_2\|_2^2 = 2$.
The variance is simply $2 \sigma^2$ (with zero overlap-dependent reduction), leading to the OTSP Signal-to-Noise Ratio:
\begin{equation}
\text{SNR}_{\text{OTSP}} = \frac{\sqrt{1-s}}{\sigma \sqrt{2}}
\end{equation}

Remarkably, we observe that:
\begin{equation}
\text{SNR}_{\text{OTSP}} = \text{SNR}_{\text{PFSR}}
\end{equation}
This is a profound, closed-form geometric result: although Löwdin orthogonalization successfully expands the clean coordinate routing margin by a factor of $\frac{1}{\sqrt{1-s}}$, it simultaneously amplifies the projection noise variance by exactly the same factor, resulting in an \textbf{exact cancellation of ensembling benefits}. Thus, the theoretical routing error probability is identical for both methods, proving that Löwdin orthogonalization provides no mathematical or empirical routing advantage under symmetric task overlap.
